\documentclass[12pt,a4paper]{article}
\usepackage[utf8]{inputenc}
\usepackage{amsmath}
\usepackage{amsfonts}
\usepackage{amssymb}
\usepackage{graphicx}
\usepackage{listings}
\usepackage{xcolor}
\usepackage{hyperref}

\title{Python Calculator Web Application}
\author{Amazon Q Workshop 2026}
\date{\today}

\lstset{
    language=Python,
    basicstyle=\ttfamily\small,
    keywordstyle=\color{blue},
    commentstyle=\color{green},
    stringstyle=\color{red},
    frame=single,
    breaklines=true
}

\begin{document}

\maketitle

\section{Introduction}
This document describes the Python Calculator web application developed as part of the Amazon Q Workshop 2026 hands-on demonstration.

\section{Project Overview}
The calculator is a Flask-based web application that provides:
\begin{itemize}
    \item Basic arithmetic operations (addition, subtraction, multiplication, division)
    \item Scientific functions (sin, cos, sqrt)
    \item Error handling for invalid operations
    \item Web-based user interface
\end{itemize}

\section{Technical Architecture}
\subsection{Backend}
\begin{itemize}
    \item \textbf{Framework:} Flask (Python web framework)
    \item \textbf{Language:} Python 3.x
    \item \textbf{Template Engine:} Jinja2
\end{itemize}

\subsection{Frontend}
\begin{itemize}
    \item \textbf{HTML5} for structure
    \item \textbf{CSS3} for styling
    \item \textbf{JavaScript} for client-side interactions
\end{itemize}

\section{Mathematical Operations}
The calculator supports the following operations:

\subsection{Basic Operations}
\begin{align}
    \text{Addition: } & a + b \\
    \text{Subtraction: } & a - b \\
    \text{Multiplication: } & a \times b \\
    \text{Division: } & \frac{a}{b}, \quad b \neq 0
\end{align}

\subsection{Scientific Functions}
\begin{align}
    \text{Sine: } & \sin(x) \\
    \text{Cosine: } & \cos(x) \\
    \text{Square Root: } & \sqrt{x}, \quad x \geq 0
\end{align}

\section{Implementation}
\begin{lstlisting}[caption=Flask Application Structure]
from flask import Flask, render_template, request, jsonify
import math

app = Flask(__name__)

@app.route('/')
def index():
    return render_template('calculator.html')

@app.route('/calculate', methods=['POST'])
def calculate():
    # Implementation details
    pass
\end{lstlisting}

\section{Deployment}
The application can be deployed using:
\begin{itemize}
    \item Local development server: \texttt{python app.py}
    \item Docker container: \texttt{docker build -t calculator .}
    \item Cloud platforms (AWS, Heroku, etc.)
\end{itemize}

\section{Testing}
Access the application at: \url{http://localhost:5000}

\section{Conclusion}
This calculator demonstrates modern web development practices using Python Flask framework, providing both educational value and practical functionality.

\end{document}
